% =====================================================================
%  note.tex —— 译注说明（凡例）
% =====================================================================
\chapter*{译注说明}
\addcontentsline{toc}{chapter}{译注说明}
\markboth{译注说明}{译注说明}

\section*{底本}

本书是 Si Li（李思，清华大学丘成桐数学科学中心）\textit{Quantum Mechanics and Geometry} 的中文译注本。
底本取自作者个人主页（\url{https://sili-math.github.io/}）公布的电子稿，封面标注为 preliminary draft、扉页标注为“updated December 2025”，共 338 页。

底本是一部面向数学系研究生的量子力学讲义：第一部分（第 1–2 章）用希尔伯特空间的公理化语言与路径积分两种形式讲基本物理原理；第二部分（第 3–4 章）转向算子层面，讲伪微分算子量子化与希尔伯特空间上算子的谱理论；第三部分（第 5–6 章）讲几何量子化与形变量子化。全书的特色是始终把“量子化”作为一个数学结构来讨论，而不是作为一套计算规则来传授。

\section*{体例}

\begin{enumerate}[leftmargin=2.2em,itemsep=2pt]
\item \textbf{页码。} 正文每页起点在源文件中以 \verb|% [原文 p.N]| 标注，$N$ 为底本的印刷页码（底本 PDF 页码 $=N+1$），供对照原书使用。
\item \textbf{陈述编号。} 定义、定理、命题、例、注等环境的编号\textbf{一律照抄底本}（形如“定义 1.1.5”“例 1.1.2”），未作重排；底本在同一节内对所有陈述环境共用单一计数器，本译注本保持这一习惯。
\item \textbf{公式。} 底本的公式\textbf{不编号}；少数由作者手标 $(*)$ 的公式保留该标签，正文中称“式 $(*)$”。
\item \textbf{图。} 底本用 TikZ 直接绘图，图均无题注、无编号，本译注本按原位置置图：多数为由原页面裁切、加注了原书页号的图；只有几何线条与英文标注的简单示意图（如双缝实验装置）按原文用 TikZ 重绘，图内标注译为中文。所有图均不加题注、不编号，以免与原书交叉引用冲突。
\item \textbf{“本节重点”框。} 每节末尾的灰蓝框是本译注本添加的教学性概述，不是底本原文，用于快速把握该节的数学内容与结论。不需要时可关闭（源文件中 \verb|\showtakeawayfalse|）。
\item \textbf{译者批注。} 以灰底小框排在被注段落之后（源文件中的 \verb|\tutornote{...}|，另有《译注汇总》一册集中收录）。批注只做三件事：标出底本的笔误与记号冲突、说明术语取舍、补出原文省略的验证步骤。批注\textbf{不改动数学内容}：凡与底本不一致处一律照录原文，另加说明。
\item \textbf{术语。} 统一依据《术语对照表》；术语以书末《英汉术语对照》为准；多数术语在中文里首次出现时以全角括号附英文原词，其余直接用中文。人名原则上译出，书末附录末节列出对照。
\end{enumerate}

\section*{免责与勘误}

翻译是对底本的一次阅读。凡译注本与底本不一致处，\textbf{以底本为准}；底本本身若有笔误，译文照原意处理并以批注标出，不擅改数学内容。
读者若发现误译、漏译或公式排印错误，欢迎指出，以便修订。

\vspace{1.2em}
\begin{flushright}
译注者\\
2026 年
\end{flushright}
